\scanpage{24}
\mspage{15}
% archive sheet number: 24
% lines 1--2
\uuline{Plus} les axiomes \inserted{des magasins} reliant \uline{$\leq$, $\ll$, $\circ$ sur $M$}, savoir Mag 1--Mag 4.

% lines 3--8
Ainsi, qu'on parte de $(\unsure{\mathscr{F}|\mathcal{F}},\leq,\ll,\mid\circ\mid)$, ou de $(\mathscr{F},M,\triangleleft,\ll,\mid\circ\mid)$, il nous faut, en plus de Mag 1--Mag 4, cinq autres axiomes.
Si on part de $(M,\leq,\ll,\mid\circ\mid,\mathscr{F})$, il nous faut un seul axiome (à tiroirs\ldots) supplémentaire, concernant le choix de $\mathscr{F}$.

% lines 9--13
Tout ceci concerne les atlas que je devrais appeler ``\unsure{spaciaux|spatiaux|spéciaux}''. Il faut maintenant ouvrir une parenthèse sur les \unsure{``atlas ``polyédraux''|``atlas'' ``polyédraux''}, qui représentent une \inserted{espèce de} structure très légèrement différente.

% lines 14--17
\uline{Définition} Le magasin $M$ est dit \inserted{pré}\uline{polyédral} si \deleted{$\forall X\in M$, $\widetilde{X}$ est un ens. ordonné polyédral, i.e.}\inserted{$(M,\leq)$ est un ens. ordonné} polyédral, i.e. si deux éléments $X,Y$ de $M$ qui sont majorés et minorés, ont un Inf \inserted{dans $M$}.

% lines 18--21: cancelled block, including its interlinear revision
\deleted{\inserted{Il est dit \uline{polyédral} si de plus, cette propriété est}}
\deleted{(\uline{NB} Si les $M_{a\leq\cdot\leq b}$ sont finis, \unsure{cette exigence|c'est l'exigence} est \unsure{entraînée|entendue} -- plus \unsure{gén.}, si $M$ est polyédral et si les $M_{a\leq\cdot\leq b}$ sont \uuline{inductifs}, alors \unsure{deux} \unsure{max. éléments}}

% lines 22--26
On introduit \deleted{\unsure{donc}} les relations plus fortes que $\ll$ (pour les magasins, pas néc. prépolyédraux)
\[
\begin{aligned}
X\underset{\mathrm{pol}}{\overset{\unsure{\mathrm{pr}|\mathrm{pré}}}{\ll}}Y
\Longleftrightarrow{}&X\ll Y,\ \text{\uuline{et}}\ \forall Y'\in\widetilde{Y},\quad
\widetilde{X}_{Y'}\overset{\text{déf}}{=}\widetilde{X}\cap\operatorname{Oubl}(\unsure{Y|Y'}
% right edge clips the Oubl argument; no closing parenthesis is supplied.
\end{aligned}
\]
est vide \uline{ou} a un plus grand élément, (i.e. $X_{Y'}$ est une figure élémentaire)

% bottom boxed margin note
\margin{En fait $X_{Y'}=X\unsure{\bullet|\cdot}Y'=\operatorname{Inf}^{\ll}(X,Y')$ dans $\mathscr{F}_M$, dans le formalisme des figures\ldots}
